\chapter{Gamification}
\label{g:gamification}

\section{Meccaniche e dinamiche (Molina/Toda)}
Gli elementi implementati, classificati secondo la tassonomia di Toda et
al.~\cite{toda}:
\begin{itemize}
\item \textbf{Performance}: punti XP, livelli (0/100/500), streak (+25/3),
soglia quiz 80\,\%, feedback immediato con explanation come hint dopo
errori ripetuti.
\item \textbf{Ecological}: vite (3), energia boss (3/turno), economia chiusa
senza valuta acquistabile.
\item \textbf{Personal}: collezione carte (deck da ricompense 1-di-3),
badge di topic e di boss, capitoli in sequenza con prerequisiti.
\item \textbf{Fictional}: narrativa dei tre boss e del dungeon-roadmap.
\item \textbf{Social}: classifica XP dell'Hub (classification
\path{overall_xp} su \path{xp}, board remota con stati vuota/offline in
app); nessuna leaderboard locale, nessuna condivisione (scelta di scope).
\end{itemize}

\section{Lente Octalysis}
Lettura secondo Chou~\cite{chou}: il nucleo \`e
\textbf{Sviluppo/Realizzazione (CD2)} nella soglia 80\,\% e nei livelli,
\textbf{Creativit\`a/Feedback (CD3)} nella scelta 1-di-3 e nel deck,
\textbf{Possesso (CD4)} nelle carte e nei badge,
\textbf{Imprevedibilit\`a (CD7)} nella pesca boss dai topic del capitolo,
\textbf{Perdita (CD8)} nel rischio boss (retry senza XP).
Significato/Narrativa (CD1) \`e coperto da intro di capitolo, lore
pre-fight ed epilogo con statistiche; Influenza sociale (CD5) vive solo
nella classifica Hub (medio, remota e best-effort); Scarsit\`a (CD6) nel
bonus giornaliero una-tantum (medio)
(visualizzazione dei pesi in Fig.~\ref{fig:octalysis}).

\begin{figure}[ht]
\centering
\begin{tikzpicture}[
  every node/.style={font=\footnotesize},
  forte/.style={draw, rounded corners=2pt, fill=green!15, align=center, inner sep=3pt},
  medio/.style={draw, rounded corners=2pt, fill=yellow!15, align=center, inner sep=3pt},
  debole/.style={draw, dashed, rounded corners=2pt, fill=gray!10, text=black!70, align=center, inner sep=3pt},
  lead/.style={thin, gray}
]
\node[regular polygon, regular polygon sides=8, minimum size=3.2cm, draw, thick, rotate=22.5] (oct) {};
% pallini peso sui vertici: verde = forte, giallo = medio, grigio = debole
\foreach \a/\c in {90/green!70!black,45/green!70!black,0/green!70!black,315/yellow!80!black,270/yellow!80!black,225/green!70!black,180/green!70!black,135/green!70!black}
  \fill[\c] (\a:1.73cm) circle (2.2pt);
% etichette esterne con leader-line corte
\node[forte, above=8mm of oct] (cd2) {CD2 $\cdot$ forte\\Sviluppo};
\node[forte, above right=6mm and 12mm of oct] (cd3) {CD3 $\cdot$ forte\\Creativit\`a};
\node[forte, right=10mm of oct] (cd4) {CD4 $\cdot$ forte\\Possesso};
\node[medio, below right=6mm and 12mm of oct] (cd5) {CD5 $\cdot$ medio\\Sociale};
\node[medio, below=8mm of oct] (cd6) {CD6 $\cdot$ medio\\Scarsit\`a};
\node[forte, below left=6mm and 12mm of oct] (cd7) {CD7 $\cdot$ forte\\Imprevedibilit\`a};
\node[forte, left=10mm of oct] (cd8) {CD8 $\cdot$ forte\\Perdita};
\node[forte, above left=6mm and 12mm of oct] (cd1) {CD1 $\cdot$ forte\\Significato};
\draw[lead] (oct.north) -- (cd2.south);
\draw[lead] (45:1.73cm) -- (cd3.south west);
\draw[lead] (0:1.73cm) -- (cd4.west);
\draw[lead] (315:1.73cm) -- (cd5.north west);
\draw[lead] (oct.south) -- (cd6.north);
\draw[lead] (225:1.73cm) -- (cd7.north east);
\draw[lead] (180:1.73cm) -- (cd8.east);
\draw[lead] (135:1.73cm) -- (cd1.south east);
\end{tikzpicture}
\caption{Ottagono Octalysis: pesi dei Core Drive nel progetto (dettagli nel testo).}
\label{fig:octalysis}
\end{figure}

\medskip
\section{Chiusura Fase 1B (voci A--E)}
\begin{itemize}
\item \textbf{A -- Narrativa (CD1)}: intro evocativa per capitolo,
mostrata una sola volta per save (\path{seenChapterIntros}); lore estesa
pre-fight per ogni boss (dialogo con riga ``Tratto'' della voce C);
epilogo ``Campagna completata!'' con statistiche della run (XP, livello,
streak massima, ricompense, boss vinti), mostrato una volta per
completamento (\path{campaignCompletionSeen}).
\item \textbf{B -- Titoli + avatar (CD4)}: un titolo per capitolo
(\textit{Sentinella della Rete}, \textit{Custode dei Dati},
\textit{Architetto del Web}), assegnato al completamento del capitolo
(tutti i topic + boss sconfitto); l'ultimo vinto resta attivo
(\path{activeTitle}). Avatar persistito: 3 icone (mago, volpe, robot)
$\times$ 3 cornici (viola, teal, arancio); i save vecchi ripartono dai
default con indici clampati.
\item \textbf{C -- Passive dei boss (CD7)}, una per boss:
Man-in-the-Middle \textit{Intercettazione} (soglia enrage al 75\,\% di HP
invece del 50\,\%: i suoi colpi diventano speciali prima, $-2$);
The Amnesiac \textit{Oblio} (ogni risposta errata costa anche $-1$
energia); Spaghetti Colossus \textit{Corazza} (assorbe il primo punto
danno da carte; i danni da quiz la aggirano).
\item \textbf{D -- Leaderboard Hub (CD5)}: classification
\path{overall_xp} (GENERAL su \path{xp}, cron lun 8:00);
\path{EngineClient.getLeaderboard()} best-effort (mai throw, lista vuota
su offline/errore) + schermata Classifica con stati di caricamento,
lista vuota e offline.
\item \textbf{E -- Daily reward (CD6)}: \textbf{+25 XP} una-tantum al
giorno (data \path{yyyy-MM-dd} in \path{lastDailyClaim}, stringa vuota =
mai riscattata; i save vecchi ripartono da claim disponibile).
\end{itemize}
